\scanpage{40}
\mspage{39}

Dans les raisonnements précédents, le fait que les éléments de $\unsure{L|\mathcal L}$ soient minimaux n'est intervenu qu'à la fin, pour prouver que la relation de \unsure{disjoint|disjoint.} $\parallel$ sur $L$ est \unsure{bien} celle induite par la disjonction des lieux, au sens de l'atlas construit. On va tout reprendre.

\noindent\begin{tabular}{@{}|p{\dimexpr\linewidth-2\tabcolsep-1pt\relax}@{}}
\uline{Définition} \deleted{\unsure{L|S}}Soit $\Lambda\subset M$. On dit que $\Lambda$ est \uline{dense} \inserted{dans $M$} si $\forall\,X,$ \deleted{\unsure{avec}} $\in M,\ \exists\,Y\in\Lambda$ avec $Y\mathrel{\overset{\circ}{\ll}}X$.
\end{tabular}

\uline{Ex} $1^{\circ})\ \Lambda=M$, exemple ``universel'' (on prend $Y=X$)

$2^{\circ})\ \Lambda=L$ \deleted{\unsure{que}} quand $\mathrm{At\,L\,1}$ est satisfait

On va énoncer des axiomes sur $(M,{\leq},{\ll},{\unsure{\ooalign{\hfil$>$\hfil\cr\hfil$<$\hfil\cr}|\gtrless}},\Lambda\subset M)$ qui assurent que $(M,{\leq},{\ll},{\ooalign{\hfil$>$\hfil\cr\hfil$<$\hfil\cr}})$ \unsure{provienne|provient} d'un atlas, et que $\Lambda\subset M$ est dense, avec \deleted{\unsure{un|en}} en plus sur l'atlas (au besoin) la condition $\mathrm{At\,D}$ ou $\mathrm{At\,D}_{\unsure{0|o}}$ ou $\mathrm{At\,C}$.

% slanted
\margin{Axiome trivial}
\begingroup\makeatletter\tagsleft@true\makeatother
\[
 X\leq Y\ \Longrightarrow\ X\ll Y\qquad\text{comme avant}
 \tag*{$\mathrm M\Lambda\,0$}
\]

% slanted
\margin{Axiome du raffinement induit}
% margin note written sideways
\margin{$\left.\rule{0pt}{3em}\right\}$ \deleted{\unsure{peu|texte} \illegible{1} \unsure{pas|pour} \unsure{de 1}}}
$\mathrm M\Lambda\,1$
\diagram{\begin{tikzcd}[ampersand replacement=\&,row sep=0pt,column sep=3pt]
 \text{Si }\ Y\deleted{\unsure{Z|X}}\ll \& X
     \& \text{et }\deleted{\unsure{x|Y}}\ \deleted{\unsure{x|Y}}Z\mathrel{\text{\uline{$\in\Lambda$}}}\quad\text{avec}\\
 \& \rotatebox{90}{$\leq$} \&\\
 Z\ll Y,\ Z\ll X',\ \deleted{\unsure{\text{Alors},\exists}} \& X'
     \& ,\ \text{alors }\exists\,Y'\text{ avec}
\end{tikzcd}}
\[
 \begin{aligned}
 &\deleted{Z'\unsure{\leq|<}Z,}\quad
 \deleted{Z'}\unsure{Y'|Z'}\leq\deleted{Z}\unsure{Y|Z},\
 \deleted{Z'}\unsure{Y'|Z'}\ll X',\ Z\ll Y'\\
 &\text{et si }\deleted{Z}\unsure{Y|Z}\leq X,\ \unsure{\mathrm{OPS}|\mathrm{ops}}\quad
 \deleted{Z'}\unsure{Y'|Z'}\leq X'
 \end{aligned}
\]

% slanted
\margin{Axiome de \deleted{la \unsure{strate minimale}} \inserted{densité}}
\[
 \forall\,X\in M,\ \exists\,Z\in\Lambda\quad\text{tel que }\ Z\mathrel{\overset{\circ}{\ll}}X
 \tag*{$\mathrm M\Lambda\,2$}
\]

% slanted
\margin{Axiome de la strate minimale}
\[
 \begin{aligned}
 &\forall\,X\in M,\ Z\in\Lambda\quad\text{avec }\ Z\ll X,\ \exists\,X'\leq X\ \text{avec}\\
 &Z\mathrel{\overset{\circ}{\ll}}X'
 \end{aligned}
 \tag*{$\mathrm M\Lambda\,3$}
\]
\endgroup

% slanted
\margin{Pour le moment il s'agit \unsure{de|des} axiomes ne faisant pas intervenir $\ooalign{\hfil$>$\hfil\cr\hfil$<$\hfil\cr}$}
\deleted{\unsure{Il est}} Il est commode d'introduire les
\[
 \begin{aligned}
 \unsure{\mathrm{omb}|\mathrm{ends}|\mathrm{ens}}_{\Lambda}(X)
    &=\{Z\in\Lambda\mid Z\ll X\}\\
 \mathrm{omb}_{\Lambda}(X)^{\unsure{\circ|0}}
    &=\{Z\in\Lambda\mid Z\mathrel{\overset{\circ}{\ll}} X\}\,.
 \end{aligned}
\]
Les raisonnements précédents montrent la validité des deux lemmes cruciaux, que je rappelle pour mémoire :
